\chapter{Requirements}
\section{Brief}
This chapter defines the requirements for the WhatsApp-like proof-of-concept developed in this dissertation. It translates the literature review into a scoped set of stakeholders, use cases, functional requirements, and non-functional requirements that are realistic for the present prototype. The emphasis is on sender-controlled photo permissions, post-send revocation, recipient understanding of changing access, and constrained forwarding to a secondary recipient. Broader multi-party mechanisms, such as co-subject intervention, remain important to the research problem but are not treated as core functional prototype requirements at this stage.

\section{Stakeholders}

\begin{enumerate}
  \item \textbf{Photo Sender:}
  The sender shares the photo and configures its permissions. This stakeholder requires clear pre-send and post-send controls, visibility over the current state of a shared photo, and interaction patterns that allow revocation or permission change without unnecessary confusion.

  \item \textbf{Primary Recipient:}
  The primary recipient receives the sender's photo in the original chat. This stakeholder requires clear feedback about whether a photo remains visible, or whether any kind of revocation action has been taken.

  \item \textbf{Secondary Recipient:}
  The secondary recipient appears only when the primary recipient forwards a photo. This stakeholder is necessary because the prototype examines whether sender-defined permissions continue to shape access after redistribution within the platform model.
 
  \item \textbf{Co-Subject:}
  The co-subject is a person depicted in a photo. This stakeholder is necessary because the prototype examines whether co-subjects can request changes to the photo's permissions.
\end{enumerate}

\section{Prototype Scope}
The current functional prototype is intentionally narrower than the full conceptual problem described in the literature review. It implements sender-side configuration of photo permissions, sender-side revision of those permissions after sending, recipient-side feedback about active or revoked access, and a constrained forwarding flow to a secondary recipient. It demonstrates these mechanisms at the interface and interaction level rather than as a production-grade messaging infrastructure. It does not implement co-subject request workflows, co-subject conflict resolution, identity or face verification, persistent audit logging, full screenshot resistance, hard guarantees against saving or offline retention, and true cross-device enforcement beyond the local proof-of-concept model.

The following concerns are treated as out of scope for the present requirements set and are better positioned in discussion and future work chapters: co-subject agency, platform-assisted removal, identity or face verification, screenshot prevention, persistent audit logging, cross-device synchronisation, and guarantees of universal deletion after disclosure.


\section{Use Cases}

\begin{enumerate}
  \item \textbf{UC-S0: Sender Configures Permissions Before Sending:}
  The sender selects or uploads a photo and decides how it should be shared before it enters the chat.

  \item \textbf{UC-S1: Sender Revises Permissions After Sending:}
  The sender revisits a previously shared photo and changes its visibility, forwarding permission, or view-once setting.

  \item \textbf{UC-R1: Primary Recipient Encounters Active, Revoked, or View-Once Access:}
  The primary recipient interprets the sender's current permission state through the chat interface and acts within those limits.

  \item \textbf{UC-R2: Primary Recipient Forwards an Allowed Photo:}
  The primary recipient forwards a photo only when redistribution is permitted by the sender's current settings.

  \item \textbf{UC-SR1: Secondary Recipient Receives Constrained Forwarded Content:}
  The secondary recipient receives a forwarded photo in a separate chat view and experiences the inherited constraints defined by the original sender.
\end{enumerate}

\clearpage
\section{Requirements}
\Needspace{12\baselineskip}
\subsection{Functional Requirements}

For clarity, the functional requirements are divided into requirements already represented in the current prototype, requirements that are only partially demonstrated or would require minor extension, and broader concerns that are explicitly deferred from the present proof-of-concept.

\subsubsection{Implemented Requirements}

\renewcommand{\arraystretch}{1.2}

\begin{longtable}{|c|P{3.8cm}|P{3.8cm}|P{2.8cm}|P{2.2cm}|}

\caption{Implemented Functional Requirements}
\label{tab:functional-req-implemented} \\

\hline
\textbf{ID} & \textbf{Requirement} & \textbf{Purpose} & \textbf{Literature Support} & \textbf{Research Gap Addressed} \\
\hline
\endfirsthead

\hline
\textbf{ID} & \textbf{Requirement} & \textbf{Purpose} & \textbf{Literature Support} & \textbf{Research Gap Addressed} \\
\hline
\endhead

\hline
\caption*{\textbf{Table~\thetable:} Implemented Functional Requirements}
\endlastfoot

\multicolumn{5}{|l|}{\textbf{1. Pre-send consent configuration}} \\ \hline
FR1 &
The prototype must allow the sender to select a photo and configure permissions before sending. &
Provide a pre-send composer that lets the sender set photo state, forwarding, view-once behaviour, expiry, and saving/export permissions before sending. &
Dynamic consent; in-context permission at send. \cite{kaye_dynamic_2015,wilbanks_design_2018,cherubini_2021,schaub_notice_2015,im_yes_2021,whatsapp_view_once,whatsapp_disappearing_messages,whatsapp_advanced_chat_privacy} &
RG1 
\\ \hline

FR6 &
The prototype must allow the sender to choose time limits, view-once, or valid-until-revoked permissions. &
Align with dynamic, revocable, and ephemeral consent as discussed in the literature. &
Consent as revertible and specific; limit duration or openness of permission; dynamic consent and temporality. \cite{im_yes_2021,robol_consent_2023,kaye_dynamic_2015,schuler_scott_dynamic_consent_2022,lee_identifying_2023,heath_2023} &
RG1
\\ \hline

FR9 &
The prototype must explicitly address screenshot and capture risk (e.g.\ warn before sharing, indicate that screenshots cannot be fully prevented). &
Treat screenshot risk as central even where perfect prevention is impossible. &
Treat screenshot risk not as marginal loophole but as central challenge; honestly communicate scope; warnings that prompt reconsideration. \cite{ingber_2025,heath_2023,schaub_notice_2015,cherubini_2021,whatsapp_disappearing_messages,whatsapp_view_once,whatsapp_advanced_chat_privacy,whatsapp_privacy_policy_uk_2025} &
RG2
\\ \hline

\multicolumn{5}{|l|}{\textbf{2. Post-send revocation and permission management}} \\ \hline
FR2 &
The prototype must allow the sender to modify or revoke access to a previously shared photo via contextual in-chat, photo-level controls (e.g.\ action menu with reply, forward, save/export, and a revocation-permissions submenu). &
Provide post-send controls via a photo-level action menu so the sender can change visibility, forwarding permission, view-once behaviour, expiry, and saving/export for an existing photo message. &
Controls that remain reachable after sending; permission revisable over time; withdrawal support. \cite{schaub_notice_2015,agrafiotis_consent_revocation_2012,merlec_smart_2021,albesher_2021,politou_2018,schnitzler_2020} &
RG1 and RG2
\\ \hline

FR3 &
The prototype must show all users clear states for active, revoked, view-once access and expiry duration. &
Render recipient-facing feedback that distinguishes visible content, revoked content, and single-view media within the chat interface. &
Users do not always understand what revocation does; clearer cues improve comprehension; legible consequences and visible outcomes. \cite{schnitzler_2020,im_yes_2021,wilbanks_design_2018,albesher_2021,schaub_notice_2015,whatsapp_delete_messages,whatsapp_disappearing_messages,whatsapp_view_once} &
RG2
\\ \hline

FR4 &
The prototype must maintain a live consent state per shared photo recording current permission status, expiry, active or revoked access, and applicable restrictions. &
Provide a usable basis for later revocation and review; without this, revocation or review is not feasible. &
Permissions as reviewable and updateable over time; explicit consent choices and consent history; revocation as closing access interval. \cite{merlec_smart_2021,robol_consent_2023,kaye_dynamic_2015,wilbanks_design_2018}&
RG1
\\ \hline

FR12 &
The prototype should provide clear confirmation or summary feedback after key consent-related actions. &
Surface the result of sending, forwarding, or changing permissions through explicit summaries or lightweight confirmation cues. &
Visible, interpretable feedback after privacy actions; clearer explanations reduce ambiguity about what changed. \cite{wilbanks_design_2018,schaub_notice_2015,schnitzler_2020,albesher_2021} &
RG1 and RG2
 \\ \hline

FR14 &
The prototype should allow the sender to attach an optional explanation when changing permissions. &
Provide an optional, non-mandatory rationale field that supports revocation transparency without requiring disclosure of sensitive reasons. &
Optional explanation supports socially intelligible boundary renegotiation; improve transparency without forcing disclosure of sensitive reasons.\cite{petronio_cpm_2020,daft_lengel_mrt_1986,short_williams_christie_1976,ishii_etal_2019,wilbanks_design_2018,schaub_notice_2015} &
RG3
\\ \hline

\multicolumn{5}{|l|}{\textbf{3. Redistribution and forwarding control}} \\ \hline
FR7 &
The prototype must constrain forwarding according to the sender's current permission state. &
Allow forwarding only when the sender has enabled redistribution and block the forwarding path otherwise. &
Forwarding changes information flow; redistribution should follow sender-defined transmission constraints; downstream sharing is a central privacy risk.\cite{nissenbaum_contextual_2011,cherubini_2021,zhang_go-sharing_2023,liu_privacy_2023,whatsapp_advanced_chat_privacy,whatsapp_advanced_chat_privacy_blog_2025,heath_2023,ingber_2025} &
RG2
\\ \hline

FR8 &
The prototype must propagate sender-defined permissions to forwarded content within the proof-of-concept model. &
Ensure that a forwarded photo shown to the secondary recipient inherits the original sender's visibility, forwarding, and view-once constraints as represented by the interface. &
Preserve original transmission principles across forwarding; inherited constraints maintain context and reduce downstream privacy loss.\cite{nissenbaum_contextual_2011,cherubini_2021,zhang_go-sharing_2023,liu_privacy_2023,whatsapp_advanced_chat_privacy,whatsapp_advanced_chat_privacy_blog_2025,heath_2023} &
RG2 and RG3
\\ \hline

FR13 &
The prototype should make the forwarded-recipient state more explicit. &
Use clearer forwarded-content labels or explanatory text so that the third-phone flow is immediately interpretable during evaluation. &
Make redistribution state legible; forwarding changes audience and context, so secondary-recipient status should be explicitly marked. \cite{nissenbaum_contextual_2011,cherubini_2021,zhang_go-sharing_2023,liu_privacy_2023,schaub_notice_2015,wilbanks_design_2018,albesher_2021,schnitzler_2020} &
RG2
\\ \hline

\multicolumn{5}{|l|}{\textbf{4. Multi-party and negotiated consent} (continued)} \\ \hline
FR5 &
The prototype must notify the recipient with a neutral system notice when access is revoked, expiry changes, or restrictions are updated. &
Keep the recipient informed whilst using socially careful wording, given that revocation can create interpersonal tension. &
Revocation as expectation-management; notices can generate mistrust or ambiguity; design for face-saving and minimal interpersonal blame. \cite{schnitzler_2020,petronio_cpm_2020,cherubini_2021,ingber_2025,daft_lengel_mrt_1986,short_williams_christie_1976,ishii_etal_2019,wilbanks_design_2018} &
RG3
\\ \hline

FR10 &
The prototype must support multi-party photos: let the sender mark that a photo contains other people, let depicted co-subjects request restriction or removal, and let the sender indicate shared versus sender-only content. &
Reflect the literature's emphasis on multi-party privacy conflicts rather than assuming only the sender matters. &
Multi-party privacy conflicts; depicted others with no opportunity to approve or shape disclosure; support for co-subject agency and shared stakes. \cite{cherubini_2021,liu_privacy_2023,bourgeus_vandercruysse_2024,xu_my_2017,xu_trust-based_2019,petronio_cpm_2020,nissenbaum_contextual_2011} &
RG3 and RG4
\\ \hline

FR11 &
The prototype must allow the recipient to request extended access, re-opened permission, or permission to keep the photo longer, with the sender deciding whether to approve. &
Support consent as socially negotiated rather than only technically toggled. &
Revocation as re-coordination of privacy boundaries; consent negotiated in unfolding interaction; permission revisited, clarified, renegotiated. \cite{petronio_cpm_2020,child_petronio_cmc_privacy_2011,im_yes_2021,kaye_dynamic_2015,schuler_scott_dynamic_consent_2022,lee_identifying_2023} &
RG3
\\ \hline

\end{longtable}

\Needspace{12\baselineskip}

\subsubsection{Deferred / Figma-Demonstrated Requirements for Co-Subject Support}

The following requirements were not implemented in the coded proof-of-concept because of time and technical limitations, but are represented in the Figma prototype to address the multi-party privacy concerns identified in the literature review.

\renewcommand{\arraystretch}{1.2}

\begin{longtable}{|c|P{3.8cm}|P{3.8cm}|P{2.8cm}|P{2.2cm}|}
\caption{Deferred / Figma-Demonstrated Functional Requirements for Co-Subject Support}
\label{tab:functional-req-cosubject} \\

\hline
\textbf{ID} & \textbf{Requirement} & \textbf{Purpose} & \textbf{Literature Support} & \textbf{Research Gap Addressed} \\
\hline
\endfirsthead

\hline
\textbf{ID} & \textbf{Requirement} & \textbf{Purpose} & \textbf{Literature Support} & \textbf{Research Gap Addressed} \\
\hline
\endhead

\hline
\caption*{\textbf{Table~\thetable:} Conceptual and Figma-Level Requirements for Co-Subject Support}
\endlastfoot

FGR1 &
The prototype should allow the sender to indicate that a photo contains one or more other people before sending. &
Introduce co-subject visibility at the point of disclosure and avoid assuming that only the sender has privacy stakes in the image. &
Multi-party privacy conflicts; co-owned content; depicted others may be affected by disclosure. \cite{cherubini_2021, liu_privacy_2023, nissenbaum_contextual_2011, petronio_cpm_2020} &
RG3 and RG4
\\ \hline

FGR2 &
The prototype should allow a depicted co-subject to receive a notification that they appear in a shared photo when the sender has marked the image as multi-party content. &
Make co-subject involvement visible and support awareness that a privacy-relevant disclosure has occurred. &
Co-subjects often lack opportunity to approve or shape disclosure; awareness is a prerequisite for agency. \cite{cherubini_2021, liu_privacy_2023, xu_my_2017, bourgeus_vandercruysse_2024} &
RG3
\\ \hline

FGR3 &
The prototype should allow a co-subject to request restriction, revocation, or removal of access to a shared photo. &
Represent co-subject agency within the permission lifecycle rather than treating privacy control as sender-only. &
Co-owned image privacy; negotiated privacy boundaries; depicted others may have legitimate competing interests. \cite{cherubini_2021, liu_privacy_2023, petronio_cpm_2020, nissenbaum_contextual_2011} &
RG3
\\ \hline

FGR4 &
The prototype should allow a co-subject to attach an optional explanation when making a restriction or removal request. &
Support socially intelligible negotiation without forcing disclosure of sensitive reasons. &
Boundary renegotiation is interpersonal and contextual; explanation can reduce ambiguity while preserving dignity. \cite{petronio_cpm_2020, short_williams_christie_1976, daft_lengel_mrt_1986, wilbanks_design_2018} &
RG3
\\ \hline

FGR5 &
The prototype should allow the sender to review and respond to a co-subject’s request through a dedicated decision interface. &
Demonstrate how multi-party privacy could be operationalised through negotiation rather than unilateral hidden system logic. &
Precautionary/collaborative mechanisms; privacy negotiation among stakeholders; consent as revisable and relational. \cite{cherubini_2021, xu_my_2017, xu_trust-based_2019, petronio_cpm_2020, im_yes_2021} &
RG3
\\ \hline

FGR6 &
The prototype should provide neutral outcome notices to affected parties when a co-subject request leads to changed permissions. &
Communicate changes clearly while minimising blame and interpersonal friction. &
Revocation notices are socially meaningful; design should support face-saving and expectation management. \cite{schnitzler_2020, petronio_cpm_2020, ingber_2025, wilbanks_design_2018} &
RG3
\\ \hline

FGR7 &
The prototype should distinguish between sender-owned content and shared-stakes content in the interface. &
Help users interpret why some photos are subject to additional review or multi-party controls. &
Privacy depends on actors, attributes, and transmission principles; co-owned photos require different treatment from sender-only content. \cite{nissenbaum_contextual_2011, cherubini_2021, liu_privacy_2023} &
RG3 and RG4
\\ \hline

FGR8 &
The prototype should show how co-subject-related restrictions affect forwarding and secondary-recipient access within the platform model. &
Demonstrate that co-subject privacy concerns continue after redistribution and are not limited to the original chat. &
Forwarding changes context and audience; downstream dissemination intensifies multi-party privacy risk. \cite{zhang_go-sharing_2023, liu_privacy_2023, nissenbaum_contextual_2011, cherubini_2021} &
RG3
\\ \hline

FGR9 &
The prototype should allow the interface to explain the limits of co-subject protection, including that offline capture, screenshots, or prior viewing cannot be fully undone. &
Keep co-subject controls realistic and avoid overstating the effect of platform-level revocation. &
Revocation is prospective and partial; screenshotting and retention remain structural limits. \cite{politou_2018, heath_2023, ingber_2025, schnitzler_2020} &
RG2 and RG3
\\ \hline

\end{longtable}

\Needspace{12\baselineskip}

\subsubsection{Deferred Concerns}

The following issues remain important to the overall dissertation problem, but they are not included as core requirements for the present prototype: co-subject request workflows, co-subject conflict resolution, identity or face verification, persistent audit logging, screenshot resistance, hard guarantees against saving or offline retention, and true cross-device enforcement beyond the local proof-of-concept model.


\Needspace{12\baselineskip}
\subsection{Non-Functional Requirements}

The non-functional requirements below are limited to qualities that can be defended for the current proof-of-concept. They focus on usability, clarity, consistency, and realistic communication of the limits of revocation rather than on infrastructure-level guarantees that the prototype does not implement.

\renewcommand{\arraystretch}{1.2}

\begin{longtable}{|c|P{3.8cm}|P{3.8cm}|P{2.8cm}|P{2.2cm}|}

\hline
\textbf{ID} & \textbf{Requirement} & \textbf{Purpose} & \textbf{Literature Support} & \textbf{Research Gap Addressed} \\
\hline
\endfirsthead

\hline
\textbf{ID} & \textbf{Requirement} & \textbf{Purpose} & \textbf{Literature Support} & \textbf{Research Gap Addressed} \\
\hline
\endhead

\hline
\caption*{\textbf{Table~\thetable:} Non-Functional Requirements}
\endlastfoot

\multicolumn{5}{|l|}{\textbf{Usability}} \\ \hline
NFR1 &
The interface must be easy to understand for first-time users. &
Ensure that users can understand the permission model and core controls immediately without prior training or repeated use. &
First-use interpretability; consent controls should be usable, legible, and consistently presented rather than hidden or ambiguous. \cite{schaub_notice_2015,wilbanks_design_2018,im_yes_2021,albesher_2021,schnitzler_2020,bourgeus_vandercruysse_2024,nouwens_dark_2020} &
RG1 and RG2
\\ \hline

NFR7 &
The system shall support quick completion of routine sharing tasks whilst introducing only proportionate friction for privacy-sensitive actions; enough to support reflection, but not so much that users avoid the feature. &
Balance everyday efficiency with reflective intervention so privacy-sensitive actions receive attention without making the system feel burdensome or disruptive. &
Consent interactions should be unburdensome but reflective: routine actions should stay low-friction, while privacy-sensitive sharing can justify proportionate prompts or warnings that support attention without causing fatigue or avoidance. \cite{im_yes_2021,schaub_notice_2015,wilbanks_design_2018,lindegren_evaluation_2021,cherubini_2021,lee_identifying_2023,baruh_big_2017} &
RG1
\\ \hline

NFR4 &
The interface must provide readable and distinguishable permission states, labels, and controls. &
Make permission-related information easy to read, compare, and interpret so users can recognise states and act confidently. &
Permission states and controls must be visually legible and clearly distinguishable so users can interpret privacy status changes, locate actions, and understand consequences without ambiguity. \cite{schaub_notice_2015,wilbanks_design_2018,im_yes_2021,schnitzler_2020,albesher_2021} &
RG1 and RG2
\\ \hline

\multicolumn{5}{|l|}{\textbf{Performance}} \\ \hline
NFR2 &
Permission-related actions and status changes should remain responsive enough for live demonstration in a local browser environment. &
Preserve a smooth and credible interaction flow by ensuring that permission changes appear promptly during use and evaluation. &
Low-friction interaction is important in fast-paced messaging; consent controls should remain usable at the moment of action and not introduce disruptive delay. \cite{wilbanks_design_2018,im_yes_2021,schaub_notice_2015,lindegren_evaluation_2021,lee_identifying_2023,albesher_2021} &
RG1
\\ \hline

\multicolumn{5}{|l|}{\textbf{Social Impact}} \\ \hline
NFR3 &
The interface must use neutral, non-blaming language to reduce unnecessary social tension when permissions change. &
Reduce avoidable interpersonal friction by ensuring that permission changes are communicated in a socially careful and non-accusatory way. &
Revocation is a socially meaningful boundary change, so interface wording should minimise blame, ambiguity, and interpersonal friction while keeping changes intelligible. \cite{petronio_cpm_2020,child_petronio_cmc_privacy_2011,schnitzler_2020,daft_lengel_mrt_1986,short_williams_christie_1976,ishii_etal_2019,wilbanks_design_2018,schaub_notice_2015} &
RG3
\\ \hline

\multicolumn{5}{|l|}{\textbf{Consistency}} \\ \hline
NFR5 &
Permission menus, indicators, and states should remain consistent with familiar WhatsApp-like interaction patterns. &
Improve learnability and reduce confusion by aligning new permission controls with interaction patterns users already recognise from WhatsApp-like messaging. &
Consistency supports interpretable privacy control: revocation-related settings should follow familiar messaging patterns and avoid inconsistent placement or presentation that makes actions harder to locate or understand. \cite{albesher_2021,schaub_notice_2015,wilbanks_design_2018,im_yes_2021,schnitzler_2020,whatsapp_delete_messages,whatsapp_disappearing_messages,whatsapp_view_once} &
RG1 and RG4
\\ \hline

\multicolumn{5}{|l|}{\textbf{Transparency}} \\ \hline
NFR6 &
The system must not over-promise. It should clearly indicate that revocation may not affect screenshots, saved copies, exports, or prior viewing. &
Calibrate user expectations by communicating the practical limits of revocation clearly and honestly. &
Revocation in messaging is only partial and prospective, so the interface should communicate scope honestly and avoid implying complete erasure where screenshots, saved copies, exports, or prior viewing remain outside full platform control. \cite{politou_2018,schnitzler_2020,heath_2023,ingber_2025,whatsapp_delete_messages,whatsapp_disappearing_messages,whatsapp_view_once,whatsapp_advanced_chat_privacy,whatsapp_privacy_policy_uk_2025,schaub_notice_2015} &
RG2
\\ \hline


\end{longtable}

\section{Feasible Extensions Consistent with Current Scope}
If further implementation time is available, the most suitable additions are small interface-level improvements that strengthen the current revocation story without expanding the project into a different problem domain. The strongest candidates are an optional sender explanation for permission changes, clearer confirmation feedback after revocation events, and more explicit labelling of forwarded content on the third phone. These additions would support evaluation and dissertation coherence without reintroducing the larger co-subject or infrastructure-heavy issues that are intentionally outside the current prototype scope.
